\documentclass{article}
\usepackage[T1]{fontenc}
\usepackage{lmodern}
\usepackage{graphicx}
\usepackage{amsmath,amssymb}
\usepackage[a4paper,margin=2cm]{geometry}
\usepackage[absolute,overlay]{textpos}
\setlength{\TPHorizModule}{1mm}
\setlength{\TPVertModule}{1mm}
\usepackage[indonesian]{babel}
\usepackage{hyperref}
\setlength{\emergencystretch}{2em}
% unit: Halaman 36

\begin{document}

% === Halaman 36 (llm) ===

\section*{Kriteria Steganografi yang Bagus}

\begin{enumerate}
    \item \textcolor{red}{Imperceptible}\
    Keberadaan pesan rahasia tidak dapat dipersepsi secara visual atau secara audial
    \item \textcolor{red}{Fidelity.}\
    Kualitas \textit{cover-object} tidak jauh berubah akibat penyisipan pesan rahasia.
    \item \textcolor{red}{Recovery.}\
    Pesan yang disembunyikan harus dapat diekstraksi kembali.
    \item \textcolor{red}{Capacity}\
    Ukuran pesan yang disembunyikan sedapat mungkin besar
\end{enumerate}

\noindent Catatan: \textit{Robustnes} bukan isu penting di dalam steganografi

\end{document}
